\chapter{Finansies en Groei}

\section{Geld stories}

Welkom die by die Graad 10 hoofstuk oor Finansiële Wiskunde. Ons gaan wiskundige vaardighede gebruik wat
jy heel waarskynlik nou gaan nodig kry.\par

As jy R~$1 000$, het, kan jy dit in jou beursie hou, of jy kan dit deponeer in ’n bankrekening. As dit in jou beursie bly,
kan jy dit enige tyd uitgee wanneer jy wil. As die bank daarna kyk vir jou, kan hulle dit aanwend met die doel om
wins daaruit te maak. Die bank "betaal" jou gewoonlik om geld te deponeer in ’n rekening om jou aan te moedig
om met hulle sake te doen. Hierdie betaling is soos ’n beloning, wat vir jou die rede is om die geld liewer in die
bank te los vir ’n rukkie as om die geld in jou beursie te hou.\par Ons noem hierdie beloning "rente".
As jy geld in ’n bankrekening deponeer, is jy eintlik besig om jou geld aan die bank te leen - en jy kan verwag om
rente te ontvang van die bank. Net so, as jy geld leen van ’n bank (of van ’n afdelingswinkel, of ’n motorhandelaar,
byvoorbeeld) dan kan jy verwag om rente te betaal op die lening. Dit is die prys vir die leen van geld.
.\par

Die idee is eenvoudig, en tog is dit die kern van die w\^{e}reld van finansies. Boekhouers, rekenmeesters en
bankiers, bestee hulle hele professionele loopbaan deur te werk met die gevolge van rente op
finansiële aangeleenthede.\par

\section{Enkelvoudige rente}
\Definition{Enkelvoudige rente}{Enkelvoudige rente is wanneer jy rente verdien op die aanvanklike bedrag wat jy belê het,
maar nie rente op rente nie.}
  
As ’n maklike voorbeeld van enkelvoudige rente, dink hoeveel jy sal kry deur R~$1~000$ te belê vir 1 jaar by ’n bank
wat vir jou enkelvoudige rente teen $5\%$ per jaar gee. Aan die einde van die jaar sal jy rente ontvang van:\par
\begin{align*}
    \mathrm{Interest} &= \text{R}~1~000 \times 5\%\\
    &= \text{R}~1~000 \times \frac{5}{100}\\
    &= \text{R}~1~000 \times 0,05\\
    &= \text{R}~50
\end{align*}

Dus, met ’n “aanvangsbedrag" van R~$1~000$ aan die begin van die jaar, sal jou “eindbedrag" aan die einde van die
jaar dan wees:\par 
\begin{align*}
    \mathrm{Closing~Balance} &= \mathrm{Opening~Balance + Interest}\\
    &= \text{R}~1~000 + \text{R}~50\\
    &= \text{R}~1~050
\end{align*}

Ons noem soms die aanvangsbedrag in finansiële wiskunde die \textsl{Hoofsom ("Principal amount")}, wat afgekort word as $P$ (R~$1~000$ in die voorbeeld). Die rentekoers vir die tydsinterval word gewoonlik as ’n persentasie aangedui deur $i$ ($5\%$ in die voorbeeld),  en die bedrag aan rente verdien (in terme van rand) word aangedui deur $I$ (R~$50$ in die voorbeeld).\par 

So, ons sien dat:
        
\begin{align}
    I = P \times i
\end{align}

en,

\begin{align}
    \mathrm{Eindbedrag} &= \mathrm{Aanvangsbedrag + Rente} \nonumber\\
    &= P + I \nonumber\\
    &= P + (P \times i)\nonumber\\
    &= P(1 + i)
\end{align}

\Tip{Onthou persentasie is 'n getal wat geskryf is oor 'n noemer van $100$. Met ander woorde, $12\%$ kan geskryf word as $\frac{12}{100}$ --- en dus $12\%$ ook uitgedruk word in desimale vorm as $0,12$.}


Bostaande berekeninge gee vir ons 'n goeie idee van die formule vir enkelvoudige rente, maar die voorbeeld praat van 'n belegging vir slegs een jaar. As 'n belegging of lening oor 'n langer tyd gemaak word, moet dit in aanmerking geneem word. Ons gebruik die simbool 'n' om die tyds periode aan te dui en heirdie waarde moet altyd in jare uitgedruk word.\par

Die formule wat gebruik word om enkelvoudigig rente te bereken, is:
\Tip{Jaarlikse rentekoerse beteken die koers word bereken oor ’n periode van ’n jaar, (per annum) = per jaar.}

\Identity{enkelvoudige rente}{
    \begin{eqnarray}
	\label{FG:SI}
	A &=& P (1 + i \cdot n)\\
	\text{Waar:~} \nonumber\\
	A &=& \text{eindbedrag} \nonumber\\
	P &=& \text{aanvangsbedrag} \nonumber\\
	i &=& \text{rentekoers per tydsinterval} \nonumber\\
	n &=& \text{aantal tydsintervalle} \nonumber
    \end{eqnarray}
}




\begin{wex}{Berekening van rente op 'n belegging}{
    As ek R~$1 000$ vir 3 jaar deponeer
in ’n spesiale bankrekening teen $7\%$ per jaar enkelvoudige rente gee,
hoeveel geld sal ek aan die einde van hierdie tydperk hê?}{

    \westep{Ons weet dat}
    \begin{align*}
	A &= ~?\\
	P &= 1~000\\
	i &= 0,07\\
	n &= 3
    \end{align*}
    
    \westep{Skryf die formula neer}
    \begin{align*}
	A &= P(1 + i \cdot n)
    \end{align*}

    \westep{Vervang die waardes in die formule}
    \begin{align*}
	A &= 1~000(1 + 0,07 \times 3)\\
	  &= 1~210
    \end{align*}

    \westep{Verduidelik jou aantwoord}
    Die eindbedrag nadat R1 000 vir 3 jaar belê is teen ’n rentekoers
van $7\%$ per jaar, is R~$1~210$.
    }
\end{wex}


\begin{wex}{Bereken rente op 'n lening}{
    Sarah leen R~$5~000$ by haar buurvrou teen 'n ooreenkome rentekoers van $12,5\%$ p.a. Sy sal die lening in een bedrag terugbetaal aan die einde van $2$ jaar. Hoeveel sal sy moet treug betaal?}{

    \westep{Skryf die veranderlikes neer}
    \begin{align*}
	A &= ~?\\
	P &= 5~000\\
	i &= 0,125\\
	n &= 2
    \end{align*}

    \westep{Skryf die formule neer}
    \begin{align*}
	A &= P(1 + i \cdot n)
    \end{align*}

    \westep{Vervang die waardes in die formule en bereken die oorblywende veranderlike}
    \begin{align*}
	A &= 5~000(1 + 0,125 \times 2)\\
	  &= 6~250
    \end{align*}

    \westep{Verduidelik jou aantwoord}
    Aan die einde van 2 jaar, sal Sarah haar buurvrou R~$6~250$ betaal.
    }
\end{wex}


Ons kan die enkelvoudige rente formule gebruik om ontbekende inligting te verkry. As ons byvoorbeeld 'n bedrag geld bel\^e vir 'n vaste tyd om 'n sekere doel to beruik kan ons die rentekoers bereken waareen ons die geld moet bel\^e om die doel te beruik. Of ons kan bereken vir hoe lank ons die geld moet bel^e as ons weet wat die aanvangsbedgrag en die eindbedrag en die rentekoers is. 


\begin{wex}{Bepaaling van die tydperk(periode) om 'n b}{epaalde uitkoms te bereik
    As ek  R~$30~000$ deponeer in ’n spesiale
bankrekening wat $7,5\%$ enkelvoudige rente p.a betaal, vir hoeveel jaar moet
ek hierdie bedrag belê om 'n eindbedgrag van R~$45~000$ laat realiseer?}{

    \westep{Skryf die bekende veranderlikes neer}
    \begin{align*}
	A &= 45~000\\
	P &= 30~000\\
	i &= 0,075\\
	n &= ~?
    \end{align*}

    \westep{Skryf die formule neer}
    \begin{align*}
	A &= P(1 + i \cdot n)
    \end{align*}

    \westep{Vervang die bekeunde waardes en bereken die onbekende}
    \begin{align*}
	45~000 &= 30~000(1 + 0,075 \times n)\\
	\frac{45~000}{30~000} &= 1 + 0,075 \times n\\
	\frac{45~000}{30~000} -1 &= 0,075 \times n\\
	\frac{(\frac{45~000}{30~000}) -1}{0,075} &= n\\
	n &= 6,6
    \end{align*}

    \westep{Verduidelik jou aantwoord}
    Dit sal ’n periode van 6 jaar en 8 maande neem vir R~$45~000$ om te groei tot R~$30~000$ teen ’n enkelvoudige rentekoers van $7,5\%$.
    }
\end{wex}


\Tip{Om 'n meet akkurate antwoord te kry, doen al jou sakrekenaar berekende in e\"en berekening en moenie tussenstappe afrond nie. Dit sal keer dat afrondings 'n te groot invloed het op jou finale antwoord.}


\begin{wex}{Bereken die rentekoers om 'n gevraagde mikpunt te bereik}{
    Teen watter enkelvoudige rentekoers moet ek R~$2~500$ bel\^e om R~$4~000$ na $5$ jaar?}{

    \westep{Skryf die bekende veranderlikes neer}
    \begin{align*}
	A &= 4~000\\
	P &= 2~500\\
	i &= ~?\\
	n &= 5
    \end{align*}

    \westep{Skryf die formule neer}
    \begin{align*}
	A &= P(1 + i \cdot n)
    \end{align*}

    \westep{Vervang die waardes en bereken die onbekende veranderlike}
    \begin{align*}
	4~000 &= 2~500(1 + i \times 5)\\
	\frac{4~000}{2~500} &= 1 + i \times 5\\
	\frac{4~000}{2~500} - 1&= i \times 5\\
	\frac{(\frac{4~000}{2~500}) - 1}{5} &= i\\
	i &= 0,12
    \end{align*}

    \westep{Verduidelik jou antwoord}
    'n enkelvoudige rentekoers van $12\%$ is nodig om R~$2~500$ oor $5$ jaar te laat groei tot R~$4~000$.
    }
\end{wex}


\subsection{As die periode nie in jare is nie}

Dikwels kan iemand nie geld vir 'n volle jaar bel\^e nie, of hulle wil 'n lening aangaan vir 'n koter periode. Hoe bereken ons dan die eindbedrag?\par

Onthou in elke jaar is daar $12$ maande. As ek geld vir slegs $3$ maande wil bel\^e sal $n$ gelyk wees aan $\frac{3}{12}$. As ek 'n lening wil uitneem vir $7$ maande is, $n = \frac{7}{12}$.\\


\begin{exercises}{Simple Interest Exercises}
    \begin{enumerate}[label=\textbf{\arabic*}.]
	\item ’n Bedrag van R~$3~500$ word belê in ’n spaarrekening wat enkelvoudige rente betaal teen ’n koers van $7,5\%$ p.a. Bereken die eindbedrag na 2 jaar.

	\item Bereken die enkelvoudige rente in die volgende probleme:
	\begin{enumerate}
	    \item ’n Lening van R~$300$  teen ’n koers van $8\%$ vir 1 jaar.

	    \item ’n Belegging van R~$2~250$ teen ’n koers van $12,5\%$ per jaar vir 6 jaar.
	\end{enumerate}

	\item Sally wil die aantal jare bereken wat sy R~$1~000$ moet bel\^e om te groei tot R~$2~500$. Sy word 'n enkelvoudige rentekoers van $8,2\%$ p.a aangebied. Hoe lank sal dit neem vir die geld om te groei tot R~$2~500$?

	\item Ek het ’n deposito van R~$5~000$ in die bank gemaak. Sestien jaar later was die eindbedrag van hierdie
belegging R~$18~000$ teen watter koers is die geld belê indien enkelvoudige rente bereken is?\\
    \end{enumerate}

    Find the answers with the short codes:\\
    \begin{tabularx}{\textwidth}{ XXXX }
	(1)	&	(2)	&	(3)	&	(4)\\
    \end{tabularx}
\end{exercises}


\subsection{Huurkoop}

As 'n algemene re\"el, is dit nie 'n goeie idee om goes op krediet te koop nie. Dit beteken dat jy geld leen om vir 'n aankoop te betaal en dat jy dus meer gaan betaal omdat jy rente op 'n lening moet betaal. Maar, soms is daar toenisting, soos 'n yskas, waarsonder mens moeilik kan klaarkom. Meeste mense het nie vermo\"e om sukke items kontant te koop nie, dus koop hulle dit op huurkoop.\par

'n Huurkoop ooreenkoms in 'n finansi\"ele kontrak tussen die winkel en die kli\"ent oor hoe die kli\"ent gaan betaal vir die produk. Die rente op 'n huurkoop ooreenkoms is altyd enkelvoudig en word slegs gehef op die uitstaande bedrag. In die meeste gevalle word verwag dat jy 'n deposito sal betaal voordat jy die produk kan huis toe neem. Die opgeloopte lening word bereken op grond van die tydperk waarvoor jy die lening wil uitneem. Die totale leningsbedrag word dan verdeel in maandlikse paaiemente oor die leningsperiode.


\begin{wex}{Huurkoop}{
    Troy wil graag ’n ekstra hardeskyf vir sy
skootrekenaar koop teen R~$2~500$ soos dit op die internet adverteer word.
Daar is die opsie om ’n deposito van $10\%$ van die koopprys te betaal
en dan in ’n huurkoop-ooreenkoms 24 gelyke maandelikse paaiemente
te betaal waar rente bereken sal word teen $7,5\%$ per jaar enkelvoudige
rente. Bereken wat Troy se maandelikse paaiement sal wees.}{

    \westep{’n Nuwe aanvangsbedrag is nodig, want die $10\%$ deposito is kontant betaal.}
    \begin{itemize}
	\item $10\%$ van R~$2~500$ = R~$250$\\
	\item Nuwe openingsbalans, $P$ = R~$2~500$ − R~$250$ = R~$2~250$\\
	\item Rentekoers, $i = 7,5\%$\\
	\item Periode, $n = 2$ jare
    \end{itemize}

    Ons moet die eindbedrag ($A$) bepaal en dan die maandelikse
paaiemente bereken.

    \westep{Van \ref{FG:SI} ons weet dat:}
    \begin{align*}
	    A &= P(1 + i \cdot n)
    \end{align*}

    \westep{Vervang die waardes en bereken die onbekende veranderlike}
    \begin{align*}
	A &= 2~250(1 + 0,075 \times 2)\\
	  &= 2~587,50\\
    \end{align*}

    \westep{Bereken die maandelikse terugbetalings op die huurkoop ooreenkoms}
    \begin{align*}
	\text{Maandelikse paaiement} &= \frac{2~587,50}{24}\\
			&= \text{R}~107,81
    \end{align*}

    \westep{verduidelik jou antwoord}
    Troy se maandelikse paaiment = R~$107,81$.
}
\end{wex}


Some voeg die winkel 'n maandelikse versekeringspremie by die maandelikse paaiemente, Hierdie premie verseker jou van meer tyd tussen 'n onbetaalde paaiement en moontlike beslaglegging op die aangekoopte produk.


\begin{wex}{Huurkoop met ekstra voorwaardes}{
    Cassidy is desperaat om 'n TV te koop en gevolglik besluit sy op 'n huurkoop ooreenkoms. Die TV se kontantsprys is R~$5~500$. Sy moet dit afbetaal oor $54$ maande teen 'n rentekoers van $21\%$p.a. 'n versekeringspremie van R~$12,50$ word bygevoeg. Hoeveel moet sy per maand betaal?}{

    \westep{Skryf die bekende veranderlikes neer}
    \begin{eqnarray*}
	A &=& ?\\
	P &=& 5~500\\
	i &=& 0,21\\
	n &=& \frac{54}{12} = 4,5
    \end{eqnarray*}
    \text{(Die vraag meld niks van 'm deposito, daarom aanvaar ons Cassidy het nie een betaal nie)}

    \westep{Skryf die formule neer}
    \begin{align*}
	    A &= P(1 + i \cdot n)
    \end{align*}

    \westep{Vervang die waardes en bereken die onbekende veranderlike}
    \begin{align*}
	A &= 5~500(1 + 0,21 \times 4,5)\\
	  &= 10~697,50\\
    \end{align*}

    \westep{Bereken die maandelikse terug betallings op die huurkoop ooreenkoms}
    \begin{align*}
	\text{Maandelikse betaalings} &= \frac{10~697,50}{54}\\
			&= 198,10
    \end{align*}

    \westep{Voeg die versekeringspremie by}
    \begin{align*}
	198,10 + 12,50 &= 210,60
    \end{align*}

    \westep{Verduidelik jou antwoord}
    Cassidy sal R~$210,60$ vir 54 months betaal tot haar TV afbetaal is.
}
\end{wex}


\Tip{Jy behoort te weet dat huurkoop teen 'n enkelvoudige rentekoers bereken word. Gebruik dus altyd die enkelvoudige rente formule in huurkoop berenkeninge.}


\begin{exercises}{Hire Purchase}
    \begin{enumerate}[label=\textbf{\arabic*}.]
	\item Vanessa wil 'n yskas op huurkoop koop. Die kontantprys is R~$4~500$. Sy moet 'n deposito van $15\%$ betaal en die oorblywende bedrag oor $24$ maande teen 'n rentekoers van $12\%$ p.a.
	\begin{enumerate}
	    \item Wat is die aanvangsbedrag (hoofsom)?
	    \item Wat is die eindbedrag wat sy terugbetaal?
	    \item Wat is Vanessa se maandlikse paaiemente?
	    \item Hoeveel het Vanessa in totaal vir haar yskas betaal?
	\end{enumerate}


	\item Bongani koop ’n eetkamertafel van R~$8~500$ op huurkoop. Hy moet enkelvoudige rente van $17,5\%$per jaar
betaal oor 3 jaar.
	\begin{enumerate}
	    \item Hoeveel sal Bongani in totaal betaal?
	    \item  Hoeveel rente betaal hy?
	    \item Wat is sy maandelikse paaiement?
	\end{enumerate}

	\item 'n Sitkamerstel word geadverteer op TV. Dit moet oor $36$ maande terugbetaal word teen R~$150$ 'n maand.
	\begin{enumerate}
	    \item Aanvaar daar is geen deposito nodig nie, Hoeveel sal die koper betaal vir die stel teen die tyd wat dit afbetaal is?
	    \item As die rentekoers $9\%$p.a is, wat is die konstantprys van die stel?\\
	\end{enumerate}
    \end{enumerate}

    Kry die oplossing met die kortkodes:\\
    \begin{tabularx}{\textwidth}{ XXX }
	(1)	&	(2)	&	(3)\\
    \end{tabularx}
\end{exercises}


\section{Saamgestelde Rente}

Saamgestelde rente laat ons toe om rente op rente te verdien. In enkelvoudig rente, verdien jy slegs rente op jou aanvanklike bedrag, maar in saamgestelde rente kan jy rente verdien op jou oorspronklike bedrag sowel as op die rente wat reeds verdien is.\par

Saamgestelde rente is 'n wonderlike konsep as jy geld bel\^e, maar, as jy geld leen sal jy meer terugbetaal as rente saamgesteld bereken word en nie enkelvoudig nie.


\Definition{Saamgestelde Rente}{Saamgestelde rente is die rente wat bereken word op die aanvangsbedrag en op die
opgeloopte rente}


Laat ons ’n formule ontwikkel vir saamgestelde rente op dieselfde manier as wat ons ’n formule ontwikkel het vir
enkelvoudige rente.\par

Indien ons aanvangsbedrag $P$ is en ons het ’n rentekoers van $i$ per jaar, dan sal die eindbedrag aan die einde
van die eerste jaar gelyk wees aan:
\begin{eqnarray*}
    \text{Eindbedrag na 1 jaar} = P(1 + i)
\end{eqnarray*}

Dit is dieselfde as enkelvoudige rente, want dit strek net oor een tydsinterval (’n jaar in hierdie geval). As ons die
geld dan onttrek en weer belê vir nog ’n jaar - soos wat ons in die uitgewerkte voorbeeld hierbo gedoen het - sal
die eindbedrag aan die einde van die tweede jaar as volg wees:
\begin{eqnarray*}
    \text{Eindbedrag na 2 jaar} &=& [P(1 + i)] \times (1 + i)\\
    &=& P(1 + i)^2
\end{eqnarray*}

As ons hierdie geld onttrek en weer vir nog ’n jaar belê, sal die eindbedrag wees:

\begin{eqnarray*}
    \text{Eindbedrag na 3 jaar} &=& [P(1 + i)^2] \times (1 + i)\\
    &=& P(1 + i)^3
\end{eqnarray*}

Ons kan sien dat die eksponent van die term $(1 + i)$ gelyk is aan die aantal tydsintervalle (jare in hierdie voor-
beeld.) Dus:


\Identity{Saamgestelde rente}{
    \begin{eqnarray}
	\label{FG:CI}
	A &=& P(1 + i)^n\\
	\text{Waar:~} \nonumber\\
	A &=& \text{eindbedrag} \nonumber\\
	P &=& \text{aanvangsbedrag} \nonumber\\
	i &=& \text{rentekoers per tydsinterval, skryf as disimaal} \nonumber\\
	n &=& \text{tydsduur van die lening bepaal} \nonumber
    \end{eqnarray}
}


\begin{wex}{Saamgestelde rente}{
    Mpho wil R~$30~000$ bel\^e in 'n rekening wat saamgestelde rente van $6\%$p.a aanbied. Hoeveel geld sal daar in die rekening wees aan die einde van $4$ jaar?}{
    
    \westep{Skryf bekende varanderlikes neer}
    \begin{eqnarray*}
	A &=& ?\\
	P &=& 30~000\\
	i &=& 0,06\\
	n &=& 4
    \end{eqnarray*}

    \westep{Skryf die formule neer}
    \begin{align*}
	A &= P(1 + i)^n
    \end{align*}

    \westep{Vervang die waardes en los op vir die onbekende veranderlike}
    \begin{align*}
	A &= 30~000(1 + 0,06)^4\\
	  &= 37~874,31
    \end{align*}

    \westep{Verduidelik jou antword}
    Mpho sal R~$37~874,31$ in die rekening h\^e na $4$ jaar.
    }
\end{wex}


\begin{wex}{Bereken die saamgestelde rentekoers om 'n mikpunt te bereik}{
    Charlie kry R~$5~000$ vir sy sestiende verjaarsdag, eerder as om dit te spandeer, besluit hy om dit te bel\^e sodat hy op sy agtiende verjaarsdag 'n deposito van R~$10~000$ kan neersit op 'n kar. Watter saamgestelde rentekoers behoort hy te kry om dit reg te kry?}{
    
    \westep{Skryf die bekende veranderlikes neer}
    \begin{eqnarray*}
	A &=& 10~000\\
	P &=& 5~000\\
	i &=& ?\\
	n &=& 2
    \end{eqnarray*}

    \westep{Skryf die formule neer}
    \begin{align*}
	A &= P(1 + i)^n
    \end{align*}

    \westep{Vervang die waardes en los op vir die oorblywende veranderlike}
    \begin{align*}
	10~000 &= 5~000(1 + i)^2\\
	\frac{10~000}{5~000}&= (1 +i)^2\\
	\sqrt[]{\frac{10~000}{5~000}} &= 1 + i\\
	\sqrt[]{\frac{10~000}{5~000}} - 1 &= i\\
	i &= 0,4142135624
    \end{align*}

    \westep{Verduidelik jou antwoord}
    Charlie moet 'n rekening kry wat 'n saamgestelde rentekoers van $41,42\%$ aanbeid.
    }
\end{wex}


\subsection{Die krag van saamgestelde rente}

Om te sien hoe belangrik "rente op rente" is, sal ons kyk na die verskil in die eindbedrae van geld wat teen
enkelvoudige rente belê is en geld wat teen saamgestelde rente belê is. Beskou ’n bedrag van R~$10~000$  wat jy
vir 10 jaar moet belê en aanvaar dat jy rente kan verdien teen $9\%$. per jaar. Wat sal die waarde van die belegging
wees na 10 jaar?\par

Die eindbedrag indien die geld enkelvoudige rente verdien, is:
\begin{align*}
    A &= P(1 + i \cdot n)\\
      &= 10~000(1 + 0,09 \times 10)\\
      &= R~19~000\\
\end{align*}

Die eindbedrag indien die geld saamgestelde rente verdien, is:
\begin{align*}
    A &= P(1 + i)^n\\
      &= 10~000(1 + 0,09)^10\\
      &= R~23~673,64
\end{align*}

As ons grafieke trek van die groei van die bedrag is dit makliker om te sien hoe die geld wat enkelvoudig bel\^e is reglynig groei maar eksponensie\"el groei as dir saamgestelde rente verien:
\begin{figure}[H]
    \begin{center}
	\begin{pspicture}(-5,-2)(7,8)
	    \psset{yunit=0.75,xunit=1}
	    \psgrid[subgriddiv=1,griddots=10,gridlabels=0](0,0)(10.4,12.3)
	    \psaxes[arrows=-, dx=1, Dx=1, dy=1, Dy=2000](0,0)(0,0)(10.4,12.3)
% 	    \psline[linecolor=red](0,5)(10,9.5)
	    \psplot[linecolor=red, algebraic=true]{0}{10}{5*((1+0.09*x))}
	\end{pspicture}
	\caption{Die groei van 'n bedrag in 'n enkelvoudige rente belegging oor $10$ jaar.}
	\label{FG:fig:SI10}
    \end{center}
\end{figure}

\begin{figure}[H]
    \begin{center}
	\begin{pspicture}(-5,-2)(7,8)
	    \psset{yunit=0.75,xunit=1}
	    \psgrid[subgriddiv=1,griddots=10,gridlabels=0](0,0)(10.4,12.3)
	    \psaxes[arrows=-, dx=1, Dx=1, dy=1, Dy=2000](0,0)(0,0)(10.4,12.3)
	    \psplot[linecolor=red, algebraic=true]{0}{10}{5*((1+0.09)^x)}
	\end{pspicture}
	\caption{Die groei van 'n bedrag wat teen saamgestelde rente bel\^e is oor $10$ jaar.}
	\label{FG:fig:CI10}
    \end{center}
\end{figure}

Dit is nog duideliker om die enorme verskil te sien indien die rentekoers $9\%$ bly en die tydspenode verleng word na $50$ jaar:
\begin{figure}[H]
    \begin{center}
	\begin{pspicture}(-2,-2)(7,8)
	    \psset{yunit=0.75,xunit=0.65}
	    \psgrid[subgriddiv=1,griddots=10,gridlabels=0](0,0)(20,15)
	    \psaxes[arrows=-, dx=2, Dx=5, dy=1, Dy=50000](0,0)(0,0)(20,15)
% 	    \psline[linecolor=red](0,0.2)(20,1.1)
	    \psplot[linecolor=red, algebraic=true]{0}{20}{0.2*((1+0.09*x*5/2))}
	\end{pspicture}
	\caption{Die groei van geld indien dit teen enkelvoudige rente bel\^e  word vir $50$ jaar.}
	\label{FG:fig:SI10}
    \end{center}
\end{figure}

\begin{figure}[H]
    \begin{center}
	\begin{pspicture}(-2,-2)(7,8)
	    \psset{yunit=0.75,xunit=0.65}
	    \psgrid[subgriddiv=1,griddots=10,gridlabels=0](0,0)(20,15)
	    \psaxes[arrows=-, dx=2, Dx=5, dy=1, Dy=50000](0,0)(0,0)(20,15)
	    \psplot[linecolor=red, algebraic=true]{0}{20}{0.2*((1+0.09)^(x*5/2))}
	\end{pspicture}
	\caption{Die groei van geld indien dit teen saamgestelde rente bel\^e  word vir $50$ jaar.}
	\label{FG:fig:CI10}
    \end{center}
\end{figure}

Weereens, hou in gedagte dat dit goeie nuus en slegte nuus is. As jy rente verdien op geld wat jy belê het,
sal saamgestelde rente daartoe lei dat die bedrag eksponensieel vermeerder. Maar, as jy geld geleen het, sal
daardie bedrag ook eksponensieel groei.


\begin{exercises}{Saamgestelde rente}
    \begin{enumerate}[label=\textbf{\arabic*}.]
	\item ’n Bedrag van R~$3~500$ word belê in ’n spaarrekening wat saamgestelde rente verdien teen $7,5\%$ per jaar.
Bereken die bedrag wat opgebou is in die rekening na verloop van 2 jaar.

	\item Morgan bel\^e R~$5~000$ in 'n rekening wat aan die einde van $5$ jaar 'n groot bedrag uitbetaal. As hy R~$7~500$ kry aan die einde van die periode, watter saamgestelde rentekoers het die bank hom aangebeid?

	\item Nicola wil geld belê teen  $11\%$ per jaar saamgestelde rente. Hoeveel geld (tot die naaste rand) moet hy belê
indien hy oor 5 jaar ’n bedrag van R~$100~000$ wil hê?\\
    \end{enumerate}

    Kry die oplossing met die kortkodes:\\
    \begin{tabularx}{\textwidth}{ XXX }
	(1)	&	(2)	&	(3)\\
    \end{tabularx}
\end{exercises}


\subsection{Inflasie}

Dikwels hoor mens ouer mense praat oor vervlo\"e tye toe alles goedkoper was. Jy kan dalk onthou toe jy as 'n kind dinge gekoop het en is nou geskok oor die verhogings in pryse. Daar is baie faktore wat die prysveranderings op 'n item be\"invloed en een daarvan is inflasie.\par

Inflasie is die gimiddelde persentasie waarteen pryse per jaar toeneem. Aangesien die pryse jaarlikes toeneem word dit bereken met die saamgestelde groei formule.


\begin{wex}{Berekening van toekomstige koste gebaseer op inflasie}
    {Melk kos R~$14$ vir $2$ liter. Hoeveel sal dit kos oor $4$ jaar as dir inflasiekoers $9\%$p.a. is?}{
    
    \westep{Skryf die bekende veranderlikes neer}
    \begin{eqnarray*}
	A &=& ?\\
	P &=& 14\\
	i &=& 0,09\\
	n &=& 4
    \end{eqnarray*}

    \westep{Skryf die formule neer}
    \begin{align*}
	A &= P(1 + i)^n
    \end{align*}

    \westep{Vervang die bekende waardes en los op vir die onbekende veranderlike}
    \begin{align*}
	A &= 14(1 + 0,09)^4\\
	  &= 19,76
    \end{align*}

    \westep{Verduidelik jou antword}
    Oor $4$ jaar sal melk R~$19,76$ vir 'n $2$ liter kos.
    }
\end{wex}


\begin{wex}{Berekening van historiese koste gebaseer op inflasie}
    {'n Boks sjokelade kan vandag R~$55$. Hoeveel soe dit $3$ jaar terug gekos het as die gimiddelde inflasiekoers $11\%$p.a. was?}{
    
    \westep{Skryf die onbekende veranderlikes neer}
    \begin{eqnarray*}
	A &=& 55\\
	P &=& ?\\
	i &=& 0,11\\
	n &=& 3
    \end{eqnarray*}

    \westep{Skryf die formule neer}
    \begin{align*}
	A &= P(1 + i)^n
    \end{align*}

    \westep{Vervang die bekende waardes en lod op vir die onbekende veranderlikes}
    \begin{align*}
	55 &= P(1 + 0,11)^3\\
	\frac{55}{(1 + 0,11)^3} &= P\\
	P  &= 40,22
    \end{align*}

    \westep{Verduidelike jou antwoord}
    Three years ago, a box of chocolates would have cost R~$40,22$.
    }
\end{wex}


\subsection{Bevolkingsgroei}

Het jy al bereken hoe groot is jou oupa en ouma se nageslag? Miskien het hulle net een kind gehad (jou Ma of Pa)of Miskien het hulle 'n dosyn kinders gehad. Hoeveel van jou ouers se broers en susters het toe weer kinders van hulle eie gehad? 'n familie stamboom begin met twee mense wat kinders het. Die kinders het dan weer kinders en so verder. Die familiboom neem eksponensie\"el toe omdat elkeen wat gebore word weer die vermo\"e het om deel te wees van verdere voortplanting. Daarom gebruik ons die saamgestede rente formule om bevolkingsgroei te bereken.


\begin{wex}{Bevolkingsgroei}
    {As die huidige bevolking van Johannesburg $3~888~180$ is, en de gimiddelde bevolkingsgroei in Suid-Afrika $2,1\%$p.a. is, wat sal die verwagte bevolking van Johannesburg in 10 jaar wees?}{
    
    \westep{Skryf die bekende veranderlikes neer}
    \begin{eqnarray*}
	A &=& ?\\
	P &=& 3~888~180\\
	i &=& 0,021\\
	n &=& 10
    \end{eqnarray*}

    \westep{Skryf die formule neer}
    \begin{align*}
	A &= P(1 + i)^n
    \end{align*}

    \westep{Vervang die waardes en los op vir die oorblywende}
    \begin{align*}
	A &= 3~888~180(1 + 0,021)^{10}\\
	  &= 4~786~342,614
    \end{align*}

    \westep{Verduidelik jou antwoord}
    Stadsbeplanners kan 'n bevolking van $4~786~343$ in Johannesburg verwag oor $10$ jaar.
    }
\end{wex}


\begin{exercises}{Inflation and Population}
    \begin{enumerate}[label=\textbf{\arabic*}.]
	\item Die gemiddelde inflasiekoers vir die afgelope aantal jaar is $7,3\%$ per jaar en jou water- en elektrisiteitsrekening is gemiddeld R~$1~425$ Bereken wat jy kan verwag om te betaal oor 6 jaar?

	\item Die prys van springmielies en koeldrank by die fliek is nou R~$60$. As die gemiddelde inflasie $9,2\%$ is, wat was die prys $5$ jaar gelde?

	\item 'n Klien dorpie in Ohio, VSA, ondervind 'n groot toename in geboortes. As die gemiddelde groeikoers van die bevolking $16\%$ p.a, hoeveel babas sal gebore word vir die $1600$ inwoners in die volgende $2$ jaar?\\
    \end{enumerate}

    Kry die oplossing met die kortkodes:\\
    \begin{tabularx}{\textwidth}{ XXX }
	(1)	&	(2)	&	(3)\\
    \end{tabularx}
\end{exercises}



\subsection{Buitelandse wisselkoerse}

Verskillende lande het hulle eie wisselkoerse ontwikkel oor die jare. In Engeland kos 'n big Mac van MacDonalds ongeveer £~$4$, in Suid-Afrika kos dit R~$20$ en in Noonwe\"e sal dit jou $48$~kr kos. In al drie hierdie lande sal jy dieselfde maaltyd kry, maar nie die selfde prys betaal nie. Toe die boek geskryf is was £~$1$ = R~$12,41$ en $1$~kr = R~$1,37$, Dit beteken dat die Big Macvir 'n Suid-Afrikaaner wat na Engeland reis, R~$49,64$ en wat in Norwe\"e reis R~$65,76$.\par

Wisselkoers het 'n veel groter invloed as die prys van 'n Big Mac. Die prys van olie styg wanneer die Suid-Afrikaanse Rand verswak. Dit gebeur omdatons met 'n swakker Rand minder in ander geldeenhede kan koop as vantevore.\par

Die wisselkoers versterk wanneer geld in land bel\^e word. Wanneer ons produkte koop wat in Suid-Afrika vervaardig word, bel\^e ons in die land se ekonomie en hou ons die geld in die land. Wanneer ons goedere koop wat ingevoer word van ander lande en gevolglik sal die Rand verswak. Onthou, hoe meer Suid-Afrikaanse goedere ons koop, hoe groter word die aanvraag daarvoor wat werkskepping sal bevorder. 'Local is lekker'!


\begin{wex}{Buitelandse Wisselkoerse}
    {Saba wil oorsee gaan om familie in Spanje te besoek. Sy het R~$10~000$ beskikbaar om te spandeer, Hoeveel Euros kan sy daarvoor koop as die wisselkoers tans €~$1$ = R~$10,68$?}{
    
    Sy kan €~$936,33$ koop met R~$10~000$.
\end{wex}


\begin{exercises}{Buitelandse Wisselkoerse}
    \begin{enumerate}[label=\textbf{\arabic*}.]
	\item Ek wil ’n iPod koop wat £~$100$, kos, en die huidige wisselkoers is £~$1$ = R~$14$. Ek glo die wisselkoers gaan verander na R~$12$ binne ’n maand.
	\begin{enumerate}
	    \item Hoeveel sal die Ipod speler in rand kos as ek dit nou koop?
	    \item Hoeveel sal ek spaar as die wisselkoers daal na R~$12$?
	    \item Hoeveel sal ek verloor as die wisselkoers verander na R~$15$?
	\end{enumerate}

	\item Bestudeer die volgende tabel met wisselkoerse:
	\begin{center}
	    \begin{tabular}{ |l|l|l| }
		\hline
		Land	&	Geldeenheid 	&	Wisselkoers\\ \hline
		Verenigde Koninkryk(VK) 	&	Ponde(£)	&	R~$14,13$\\ \hline
		Amerika (VSA) 	&	Dollars(\$}	&	R~7,04\\ \hline
	    \end{tabular}
	\end{center}
	
	\begin{enumerate}
	    \item In Suid-Afrika is die koste van ’n nuwe Honda Civic R~$173~400$. In Engeland kos dieselfde voertuig £~$12~200$ en in die \$~$21~900$. In watter land is die voertuig die goedkoopste as jy die pryse
omskakel na Suid-Afrikaanse rand ?

	    \item Sollie en Arinda is kelners in ’n Suid-Afrikaanse restaurant wat baie oorsese toeriste lok. Sollie kry ’n £~$6$ fooitjie van ’n toeris en Arinda kry \$~$12$.  Hoeveel Suid-Afrikaanse rand het elkeen gekry?
	\end{enumerate}
    \end{enumerate}

    Kry die oplossing met die kortkodes:\\
    \begin{tabularx}{\textwidth}{ XX }
	(1)	&	(2)\\
    \end{tabularx}
\end{exercises}


\subsection{Opsomming}

\begin{itemize}
    \item Daar is twee soorte rente: enkelvoudige rente en saamgetelde rente:\\
    
% \framebox{
    \begin{tabularx}{\textwidth}{ XX }
	enkelvoudige rente	&	saamgetelde rente\\
	$A = P (1 + i \cdot n)$	&	$A = P(1 + i)^n$\\
    \end{tabularx}
    \par
    Waar:
    \begin{eqnarray*}
% 	\text{Where:~}\\
	A &=& \text{Aanvangsbedrag}\\
	P &=& \text{Eindbedrag}\\
	i &=& \text{rentekoers per tydsinterval}\\
	n &=& \text{aantal tydsintervalle}
    \end{eqnarray*}
% }

    \item Huurkoop terugbetalings word bereken met die enkelvoudige renteformule op die oorblywende lening bedrag. Maandelikse terugnetaalings word bereken deur doe opgeloopte bedrag te deel deur die aantal maande waarbinne die lening moret terug betaal word .

    \item Bevolkingsgroei en inflasie word bereken met behulp van die saamgestelde rente formule.

    \item ’n Buitelandse wisselkoers is die prys van een geldeenheid in terme van ’n ander.
\end{itemize}


Die volgende 3 videos gee ’n opsomming oor hoe om rente te bereken. Let daarop dat hoewel die voorbeelde
met dollars gedoen word, ons die feit kan gebruik dat die dollar ’n desimale geldeenheid is net soos die rand
(ignoreer die wisselkoers). Dit is wat in die onderafdelings gedoen is.\par 

\begin{figure}[H]
    \textnormal{Khan Akademie video oor rente - 1}
    \vspace{.1in}
    \nopagebreak
    \raisebox{-5 pt}{ \includegraphics[width=0.5cm]{col11306.imgs/summary_www.png}} { (Video:  MG10030 )}
 \end{figure}

\begin{figure}[H] % horizontal\label{m39335*interest-2}
    \textnormal{Khan Akademie video oor rente  - 2}
    \vspace{.1in}
    \nopagebreak
    \raisebox{-5 pt}{ \includegraphics[width=0.5cm]{col11306.imgs/summary_www.png}} { (Video:  MG10031 )}
\end{figure}

Let Wel: Aan die einde van hierdie video word die reël van 72 genoem. Jy sal nie hierdie reël gebruik nie, maar
sal liewer die probeer-en-tref metode gebruik om die gevraagde probleem op te los.

\begin{figure}[H] % horizontal\label{m39335*interest-3}
    \textnormal{Khan Akademie video oor rente  - 3}
    \vspace{.1in}
    \nopagebreak
    \raisebox{-5 pt}{ \includegraphics[width=0.5cm]{col11306.imgs/summary_www.png}} { (Video:  MG10032 )}
\end{figure}


\begin{eocexercises}{Finansies en Groei}
    \begin{enumerate}[label=\textbf{\arabic*}.]
	\item Jy is met vakansie in Europa. Die hotel vra €~$200$ euro per nag. Hoeveel rand het jy nodig om die hotelrekening te betaal as die wisselkoers  €~$1$ = R~$9,20$?

	\item Bereken hoeveel rente jy sal verdien as jy R~$500$ vir 1 jaar belê teen die volgende rentekoerse:
	\begin{enumerate}
	    \item $6,85\%$ enkelvoudige rente.
	    \item $4,00\%$ saamgetelde rente.
	\end{enumerate}

	\item Bianca het R~$1~450$ om vir 3 jaar te belê. Bank A bied ’n spaarrekening aan wat enkelvoudige rente betaal
teen ’n koers van $11\%$ per annum. Bank B bied ’n spaarrekening aan wat saamgestelde rente betaal
teen ’n koers van $10,5\%$ per annum. Watter bank gaan vir Bianca die spaarrekening gee met die grootste
opgeloopte bedrag aan die einde van die 3 jaar?

	\item Hoeveel enkelvoudige rente is betaalbaar op ’n lening van R~$2~000$ vir ’n jaar indien die rentekoers $10\%$ per jaar is?

	\item Hoeveel saamgestelde rente is betaalbaar op ’n lening van  R~$2~000$ vir ’n jaar indien die rentekoers $10\%$ per jaar is?

	\item Bespreek:
	\begin{enumerate}
	    \item watter soort rente jy sal verkies as jy die lener is?

	    \item watter soort rente jy sal verkies as jy die bankier is?
	\end{enumerate}

	\item Bereken die saamgestelde rente vir die volgende probleme.
	\begin{enumerate}
	    \item ’n lening van R~$2~000$ vir 2 jaar teen  $5\%$ per jaar
	    \item ’n belegging van  R~$1~500$ vir 3 jaar teen $6\%$ per jaar
	    \item ’n lening van R~$800$ vir l jaar teen $16\%$per jaar
	\end{enumerate}

	\item As die wisselkoers vir ¥~$100$ = R~$6,2287$ (¥ = Yen) en 1 Australiese dollar (AUD) = R~$5,1094$, bepaal die wisselkoers
tussen die Australiese dollar en die Japanese jen.
	\item Bonnie het ’n stoof gekoop vir R~$3~750$. Na 3 jaar het sy dit klaar betaal, asook die R~$956,25$ huurkoopkoste.
Bereken die koers waarteen enkelvoudige rente bereken is.
    \end{enumerate}

    Kry die oplossing met die kortkodes:\\
    \begin{tabularx}{\textwidth}{ XXXXXXXXX }
	(1) & (2) & (3) & (4) & (5) & (6) & (7) & (8) & (9)\\
    \end{tabularx}
\end{eocexercises}
